\documentclass[11pt]{article}
\usepackage[a4paper,margin=2.3cm]{geometry}
\usepackage{amsmath,mathtools,amssymb,amsthm,booktabs,microtype,xcolor}
\usepackage[colorlinks,linkcolor=blue!50!black,citecolor=blue!50!black]{hyperref}
\newtheorem{theorem}{Theorem}
\newtheorem{proposition}[theorem]{Proposition}
\newtheorem{corollary}[theorem]{Corollary}
\theoremstyle{definition}
\newtheorem{definition}[theorem]{Definition}
\newtheorem{computation}[theorem]{Computation}
\theoremstyle{remark}
\newtheorem{remark}[theorem]{Remark}
\newcommand{\E}{\mathbb{E}}
\newcommand{\Var}{\operatorname{Var}}
\newcommand{\Cov}{\operatorname{Cov}}
\newcommand{\KL}{\operatorname{KL}}
\newcommand{\code}[1]{\texttt{\detokenize{#1}}}
\title{Identifiability of B's row from endpoint counts:\\ what can be proved, and what the E0 experiments actually show}
\author{BIO-SBI --- note accompanying \code{report_sequential}, Part I}
\date{2026-11-07}

\begin{document}
\maketitle

\begin{abstract}
The sequential report (Part I: E0, X5a, X5c, E0 at 4096 clones) concludes that, with B divisions free, B's row is \emph{structurally} non-identifiable from endpoint counts. We study the Galton--Watson (GW) model formally and with the exact likelihood (probability generating functions, no simulation). The results split the claim in two. (i) \emph{What is true and provable}: B's row enters the mean dynamics only through its mean offspring vector, so first moments are blind to a 3-dimensional family that trades switching for division (Theorem~\ref{thm:mean}); within that family the row is fixed by its second factorial moments, and B divisions always inflate clone-size variance (Theorem~\ref{thm:second}); and the variance of the Kesten--Stigum limit sees the second moments of \emph{both} rows through a single scalar, so A's row can absorb B's (Theorem~\ref{thm:varw}). (ii) \emph{What is not true}: the model is not structurally non-identifiable. The exact Fisher information at the truth is positive definite at every endpoint $n=2,\dots,12$ (Computation~\ref{comp:fim}), so the parameters are locally identifiable. The non-identifiability is \emph{practical}. At the real-like design the Cram\'er--Rao bound on B$\to$A is 0.28, larger than the prior sd. Globally, the best model in which B divides with probability 0.5 is only 0.36 nats worse than the truth over all 266 organoids (Computation~\ref{comp:profile}). At 4096 clones on three endpoints the bound is 0.023, and the exact likelihood rejects the NPE answer by 116--178 nats on the very datasets used in the report; the \code{moments} summaries still separate it (97\% accuracy). That failure is in the amortised inference, not in the model. Finally, the information about B's row decays geometrically with the endpoint generation: one endpoint at generation 2--3 gives 10--15 times the information of the real-like design with the same number of organoids.
\end{abstract}

\tableofcontents

%======================================================================
\section{Setting}\label{sec:setting}
Types A (= LGR5$^+$, index 1) and B (= KRT20$^+$, index 2). Each living cell has one event per generation, drawn from its type's \emph{row}: death, persist, switch, or division into AA, AB, BB. Write the rows as
\[
\pi_A=(d_A,p_A,t_A,a_A,b_A,c_A),\qquad \pi_B=(d_B,p_B,q,a_B,b_B,c_B),
\]
in the order death, persist, switch (A$\to$B resp.\ B$\to$A), div\_AA, div\_AB, div\_BB. With $s=(x,y)$ marking A and B, the offspring probability generating functions (pgfs) are
\begin{align*}
f_A(x,y)&=d_A+p_Ax+t_Ay+a_Ax^2+b_Axy+c_Ay^2,\\
f_B(x,y)&=d_B+p_By+q\,x+a_Bx^2+b_Bxy+c_By^2 .
\end{align*}
$Z_n=(Z_{n,A},Z_{n,B})$ is the clone at generation $n$ from one A founder; its pgf is $F^A_n$ with $F_n=f\circ F_{n-1}$, $F_0=\mathrm{id}$. An endpoint $n$ contributes $N_n$ i.i.d.\ clones from the law of $Z_n$ \emph{conditioned on $Z_n\neq 0$} (detected clones, perfect capture). A design is $\mathcal D=\{(n,N_n)\}$. The free parameter is $\theta=(\pi_A,\pi_B)\in\Delta_5\times\Delta_5$ (dimension 10; \code{two_type}). Fixing $a_B=b_B=c_B=0$ gives \code{reversible_ab} (dimension 7).

The truth used throughout E0/X5/E4 is
\[
\theta^\star:\ \pi_A=(.03,.03,.05,.55,.05,.29),\quad \pi_B=(.10,.10,.80,0,0,0),
\]
which lies on the boundary of $\Delta_5$ for B (no B divisions). Mean matrix $M_{kj}=\partial_j f_k(1,1)$, Perron root $\rho$, right/left Perron vectors $r,u$ ($Mr=\rho r$, $uM=\rho u$, $u\cdot r=1$). At $\theta^\star$: $M=\begin{psmallmatrix}1.18&0.68\\0.80&0.10\end{psmallmatrix}$, $\rho=1.554$, $\lambda_2=-0.274$, $r\propto(1,0.550)$.

\paragraph{Three notions of identifiability.} We separate them because the report's word ``structural'' conflates them.
\begin{definition}
(a) $\theta$ is \emph{structurally (globally) identifiable} from design $\mathcal D$ if $\theta\mapsto (\mathcal L(Z_n\mid Z_n\neq0))_{n\in\mathcal D}$ is injective. (b) It is \emph{locally identifiable} at $\theta^\star$ if this holds in a neighbourhood; a sufficient condition is a positive definite Fisher information $I_{\mathcal D}(\theta^\star)=\sum_n N_n I_n(\theta^\star)$. (c) A parameter is \emph{practically non-identifiable} at design $\mathcal D$ if it is locally identifiable but the attainable precision, e.g.\ the Cram\'er--Rao (CR) bound $\sqrt{g^\top I_{\mathcal D}^{-1}g}$ for a functional with gradient $g$, is no better than the prior spread.
\end{definition}
A model that is structurally non-identifiable stays so for every $N$. A practically non-identifiable one becomes identifiable as $N$ grows, at rate $N^{-1/2}$ (Bernstein--von Mises), and better designs move it.

%======================================================================
\section{Exact results on moments}\label{sec:moments}

\begin{theorem}[Mean fibres]\label{thm:mean}
For every $n$, $\E Z_n=e_A M^n$, and B's row enters $M$ only through its mean offspring vector
\[
\mu_B=(M_{BA},M_{BB})=(q+2a_B+b_B,\ p_B+b_B+2c_B).
\]
For fixed $\mu_B$ the admissible B rows form a polytope of dimension 3 (when $\mu_B$ is interior), spanned locally by the moves
\begin{align*}
e_1&:\ q\to q-2\varepsilon,\ a_B\to a_B+\varepsilon,\ d_B\to d_B+\varepsilon &&\text{(switch $\to$ divide-into-AA and die)},\\
e_2&:\ p_B\to p_B-2\varepsilon,\ c_B\to c_B+\varepsilon,\ d_B\to d_B+\varepsilon &&\text{(persist $\to$ divide-into-BB and die)},\\
e_3&:\ q\to q-\varepsilon,\ p_B\to p_B-\varepsilon,\ b_B\to b_B+\varepsilon &&\text{(switch and persist $\to$ divide-into-AB)}.
\end{align*}
Every expected count at every endpoint, and therefore every statistic that depends on the data only through mean counts, is constant on these fibres.
\end{theorem}
\begin{proof}
$\E[Z_{n}\mid Z_{n-1}]=Z_{n-1}M$ gives $\E Z_n=e_AM^n$. The two entries of $\mu_B$ are linear in $\pi_B$, and $\pi_B$ lies in a 5-dimensional simplex. So the fibre is the intersection of the simplex with an affine subspace of codimension 2, which has dimension 3 at interior points. Each $e_i$ keeps $\sum\pi_B=1$ and both entries of $\mu_B$; they are linearly independent.
\end{proof}

\noindent At the truth only the directions with $\varepsilon\ge0$ are admissible (B divisions $\ge0$). The C2ST alternative of E0 (B$\to$AA $=a$, B$\to$A $=0.8-2a$, B death $=0.1+a$) is exactly $e_1$.

\begin{theorem}[Second moments fix the row; B divisions inflate variance]\label{thm:second}
Let $G_k=\E[\xi_k\xi_k^\top]-\operatorname{diag}(\mu_k)$ be the second factorial moment matrix of type $k$'s offspring $\xi_k$. Then
\[
G_k=\begin{pmatrix}2a_k&b_k\\ b_k&2c_k\end{pmatrix}\qquad(k=A,B),
\]
and $\pi_k\mapsto(\mu_k,G_k)$ is an affine bijection from $\Delta_5$ onto its image. Moreover, with $V_k=\Cov(\xi_k)$ and $m_j=e_AM^j$,
\begin{equation}\label{eq:cov}
\Cov(Z_n)=\sum_{j=0}^{n-1}(M^\top)^{n-1-j}\Big(\sum_k m_j(k)V_k\Big)M^{n-1-j}.
\end{equation}
Along a mean fibre (Theorem~\ref{thm:mean}) only $V_B=G_B+\operatorname{diag}(\mu_B)-\mu_B\mu_B^\top$ changes, by $\Delta G_B$. Starting from a row without B divisions (as $\theta^\star$), every admissible move has $\Delta G_B=\begin{psmallmatrix}2a_B&b_B\\b_B&2c_B\end{psmallmatrix}$ with $a_B,b_B,c_B\ge0$. Hence, for every $w\in\mathbb R^2_{>0}$ (e.g.\ $w=(1,1)$, total clone size) and $n\ge2$,
\[
\Var(w\cdot Z_n)\ \text{strictly increases with the B-division mass}\ \delta=a_B+b_B+c_B .
\]
\end{theorem}
\begin{proof}
The formula for $G_k$ is a direct computation over the six outcomes: only the divisions have two daughters, and $\E[\xi_A(\xi_A-1)]=2a_k$, $\E[\xi_A\xi_B]=b_k$, $\E[\xi_B(\xi_B-1)]=2c_k$. Given $(\mu_k,G_k)$, the three division probabilities are read off $G_k$, then persist and switch off $\mu_k$, and death is the remainder; so the map is an affine bijection. Formula \eqref{eq:cov} is the standard covariance recursion $\Cov(Z_{n})=M^\top\Cov(Z_{n-1})M+\sum_k\E Z_{n-1,k}V_k$ (Harris 1963; Athreya and Ney 1972, Ch.~V), unrolled. Along a fibre $M$ and $\mu_B$ are fixed, so the change is
\[
\Delta\Var(w\cdot Z_n)=\sum_{j=1}^{n-1}m_j(B)\,v_j^\top\Delta G_Bv_j,\qquad v_j=M^{n-1-j}w>0,
\]
and $v^\top\Delta G_Bv=2a_Bv_A^2+2b_Bv_Av_B+2c_Bv_B^2>0$ whenever $v>0$ and $\delta>0$. Since $t_A+b_A+c_A>0$ we have $m_j(B)>0$ for $j\ge1$.
\end{proof}

\noindent Theorems~\ref{thm:mean}--\ref{thm:second} explain why the E0 C2ST found the switching/division difference only in clone-size \emph{variability}. They also explain why \code{CohortMeans} cannot work. Its 6 features for 10 parameters are rank deficient by construction, and E0 showed that variance carries the information.

\begin{theorem}[The variance of the growth limit sees one scalar]\label{thm:varw}
Let $\rho>1$ and $W=\lim_n r\cdot Z_n/\rho^n$ (a.s.\ and in $L^2$, since the offspring have finite variance). Then
\[
\Var W=\sum_{k\in\{A,B\}}\beta_k\, r^\top V_k r,\qquad \beta_k=\sum_{j\ge0}\rho^{-2(j+1)}m_j(k).
\]
Hence, on the product of the two mean fibres (both rows moved with $M$ fixed), $\Var W$ depends on the six second factorial moments $(G_A,G_B)$ only through the scalar
\[
\beta_A\,r^\top G_Ar+\beta_B\,r^\top G_Br .
\]
An increase of B's division variance can be cancelled exactly by fewer A divisions. At $\theta^\star$, $\beta=(1.000,\,0.294)$.
\end{theorem}
\begin{proof}
$W_n=r\cdot Z_n/\rho^n$ is a martingale. Its increments $W_{j+1}-W_j=\rho^{-(j+1)}\sum_{\text{cells }c\text{ in gen }j}(r\cdot\xi_c-\rho r_{k(c)})$ are uncorrelated with conditional variance $\rho^{-2(j+1)}\sum_kZ_{j,k}\,r^\top V_kr$. Take expectations and sum (orthogonality of martingale increments). $W_0=r_A$ is deterministic. Along fibres $M$, hence $\rho$, $r$, $m_j$ and $\beta$, are fixed, and $r^\top V_kr=r^\top G_kr+\sum_i\mu_k(i)r_i^2-(\mu_k\cdot r)^2$ changes only through $r^\top G_kr$.
\end{proof}

\noindent Check against the exact law (\code{scripts/b_row_theory/varW.py}): $\Var(r\cdot Z_n/\rho^n)=0.3263$ ($n=10$), 0.3286 ($n=12$) against the formula $\Var W=0.3302$.

\begin{remark}[What Theorem~\ref{thm:varw} does and does not explain]
For large clones $Z_n\approx\rho^nWu$. The dominant second-order feature of late endpoints, the spread of clone sizes, therefore constrains A's and B's division variances only jointly. This is the mechanism by which leaving B divisions free biases A's row (E0 at 4096: A persist 0.03 $\to$ 0.16, A$\to$BB 0.29 $\to$ 0.21, \S\ref{sec:e0large}). It is \emph{not} a non-identifiability. The type composition, the early endpoints and the higher moments carry further constraints. A naive compensating move that matches $M$ and $\Var W$ exactly (A: div\_AA and div\_BB $-0.22$, div\_AB $+0.44$) is \emph{easier} to detect than the uncompensated B move: KL 0.17--0.28 vs 0.011--0.016 nats per clone. The cheapest compensations are found numerically in \S\ref{sec:profile}. They raise A \emph{death}, not A persist: B$\to$AA instead of B$\to$A makes extra A cells, and A death removes them. The NPE answer at 4096 (A persist 0.15) is a suboptimal version of the same compensation.
\end{remark}

%======================================================================
\section{Exact likelihood: local identifiability and its strength}\label{sec:fim}
The law of $Z_n$ is computed exactly by evaluating $F^A_n$ on an $L\times L$ grid of roots of unity and applying a 2-D FFT. The polynomial $F_n^A$ has degree $\le 2^n$ in each variable, so $L=2^n+1$ is alias-free. For $n=10$, $L=384$ already reproduces the alias-free result to $10^{-17}$, because the mass beyond 384 cells is below $10^{-14}$. Derivatives $\partial_\theta F_n$ are propagated by forward-mode differentiation through the iteration $F_n=f(F_{n-1})$ and checked against finite differences (relative error $10^{-11}$). Means agree with $e_AM^n$ to machine precision. The Fisher information of the \emph{conditioned} law $\tilde p=p/(1-p_{00})$ is $I_n=\sum_{z\neq0}\tilde p_z\,\nabla\log\tilde p_z\nabla\log\tilde p_z^\top$, in the 10 coordinates ``entry minus death'' of each row. Because $\theta^\star$ is on the boundary, B-division directions are one-sided. A positive definite $I_n$ still gives $\KL(\theta^\star\|\theta)\ge c\,|\theta-\theta^\star|^2$ on the admissible cone near $\theta^\star$. Code: \code{scripts/b_row_theory/}.

\begin{computation}[Local identifiability at $\theta^\star$]\label{comp:fim}
For every single endpoint $n\in\{2,\dots,10,12\}$ the per-clone Fisher information $I_n(\theta^\star)$ of the 10-parameter model is positive definite (Table~\ref{tab:fim}). Hence $\theta^\star$ is locally identifiable from \emph{any single} endpoint distribution. The smallest eigenvalue is far above floating-point error: its ratio to the largest is $\ge5\times10^{-7}$, against a relative accuracy of $\approx10^{-12}$ for the computed probabilities.
\end{computation}

\begin{table}[h]\centering\small
\begin{tabular}{@{}rrrrrr@{}}\toprule
$n$ & $\lambda_{\min}(I_n)$ & cond.\ $I_n$ & CR var.\ B$\to$A & CR var.\ $\delta$ & $N$ for sd 0.05 \\\midrule
2 & 0.362 & $1.3\times10^{2}$ & 1.8 & 1.7 & 730 \\
3 & 0.331 & $1.9\times10^{2}$ & 2.3 & 1.8 & 920 \\
4 & 0.205 & $5.1\times10^{2}$ & 3.6 & 3.6 & 1\,440 \\
5 & 0.051 & $3.2\times10^{3}$ & 13 & 14 & 5\,300 \\
6 & 0.037 & $6.6\times10^{3}$ & 18 & 19 & 7\,200 \\
7 & 0.015 & $2.3\times10^{4}$ & 41 & 39 & 16\,000 \\
8 & 0.0090 & $5.3\times10^{4}$ & 64 & 54 & 26\,000 \\
9 & 0.0044 & $1.4\times10^{5}$ & 124 & 91 & 50\,000 \\
10 & 0.0023 & $3.5\times10^{5}$ & 218 & 143 & 87\,000 \\
12 & 0.00067 & $1.8\times10^{6}$ & 686 & 413 & 274\,000 \\\bottomrule
\end{tabular}
\caption{Exact per-clone Fisher information at $\theta^\star$, single endpoint, 10 free parameters. CR variance of B$\to$A and of the B-division mass $\delta=a_B+b_B+c_B$ with all other parameters free (profiled), per clone: divide by $N$. Last column: $N$ for a CR sd of 0.05 on B$\to$A from that endpoint alone.}\label{tab:fim}
\end{table}

\paragraph{The weak direction.} The eigenvector of the smallest eigenvalue of $I_{\mathcal D}$ is, at every design checked, dominated by B$\to$A ($\approx0.8$) against B$\to$AA and B$\to$AB ($-0.2$ to $-0.5$), with a smaller B persist component. This is the switching-versus-division direction of Theorem~\ref{thm:mean} (move $e_1$, partly $e_3$), now as a property of the full likelihood rather than of the means.

\begin{proposition}[Practical non-identifiability at realistic designs]\label{prop:cr}
CR bounds for B$\to$A and for $\delta$ (all other parameters free) are listed in Table~\ref{tab:cr}. At the real-like and E0 designs the bound on B$\to$A is 0.27--0.28. The uniform Dirichlet prior has sd 0.14 (Beta(1,5)) and the flat prior on B$\to$A has sd 0.29. So no estimator that is unbiased near $\theta^\star$ can be more precise than the prior, and the posterior is necessarily prior-dominated. This is the regime of E0, X5a and X5c, and their conclusions hold there.
\end{proposition}

\begin{table}[h]\centering\small
\begin{tabular}{@{}lrrr@{}}\toprule
design (endpoint: clones) & organoids & CR sd B$\to$A & CR sd $\delta$ \\\midrule
real-like $\{6{:}165,10{:}101\}$ & 266 & 0.28 & 0.28 \\
E0 $\{4,8,12\}\times32$ & 96 & 0.27 & 0.28 \\
$\{4,6,10\}\times512$ & 1\,536 & 0.065 & 0.068 \\
E0 at 4096: $\{4,6,10\}\times4096$ & 12\,288 & 0.023 & 0.024 \\\midrule
$\{4\}\times266$ & 266 & 0.12 & 0.12 \\
$\{3\}\times266$ & 266 & 0.093 & 0.083 \\
$\{2,3\}\times133$ & 266 & 0.077 & 0.066 \\
$\{3,6,10\}\times\{89,89,88\}$ & 266 & 0.12 & 0.12 \\
$\{2,4,6,10\}\times\approx66$ & 266 & 0.093 & 0.091 \\
$\{3,4\}\times1000$ & 2\,000 & 0.036 & 0.033 \\\bottomrule
\end{tabular}
\caption{Cram\'er--Rao sd at $\theta^\star$ (exact Fisher information, GW, generations). Top: designs used in the report. Bottom: same 266 organoids redistributed to early endpoints.}\label{tab:cr}
\end{table}

\begin{proposition}[The information about B's row is in early generations]\label{prop:decay}
Per clone, the information about B$\to$A and $\delta$ decreases geometrically with the endpoint. From $n=5$ to $n=12$ the CR variance of B$\to$A grows by a factor 1.4--2.3 per generation (geometric mean 1.76; Table~\ref{tab:fim}), and the condition number of $I_n$ by a factor 2--3.5. Between $n=2$ and $n=4$ the loss is milder (factor 2 in total).
\end{proposition}
\noindent \emph{Interpretation.} At small $n$ the counts are small integers, and individual events are visible in the law. A clone $(2,0)$ at $n=2$ arises from A$\to$B$\to$AA but not from A$\to$B$\to$A. As $n$ grows the law of $Z_n$ approaches $\rho^nWu$ plus Gaussian composition noise. Its first moments are blind to the fibres of Theorem~\ref{thm:mean}, and its main second moment ($\Var W$) mixes the two rows (Theorem~\ref{thm:varw}). The geometric rate is a numerical observation over $n\le12$, not a proved asymptotic.

\noindent \emph{Consequence for design.} With the same 266 organoids, an endpoint at generation 2--3 gives a CR sd of 0.08--0.09 on B$\to$A, against 0.28 at $\{6,10\}$: 10--15 times the Fisher information. Adding one early endpoint to $\{6,10\}$ already helps (0.12 for $\{3,6,10\}$). Whether generation 2--3 is reachable in the wet lab (day 2--3 for a $\approx$20\,h A cycle) and whether BH behaves alike are open (\S\ref{sec:open}).

%======================================================================
\section{E0 at 4096 clones revisited}\label{sec:e0large}
The report (\S\,``E0 revisited at 4096 clones'') concludes ``structural non-identifiability'' because, at $\{4,6,10\}\times4096$ with a 20k-simulation NPE, B$\to$A stays at 0.17--0.34 (truth 0.8) and the B-division mass at its prior mean. Three exact checks contradict the ``structural'' reading for GW.

\paragraph{1. The bound.} The CR sd is 0.023 (Table~\ref{tab:cr}). NPE reported a posterior sd of 0.11 and a bias of 0.47, which is 20 CR sd.

\paragraph{2. The exact likelihood rejects the NPE answer on the report's own data.} For each run we take the NPE posterior mean $\hat\theta$, renormalise each row, and evaluate the exact conditioned log-likelihood on the stored reference data (\code{reference.npz}, $3\times4096$ clones):
\begin{center}\small
\begin{tabular}{@{}lrrrrr@{}}\toprule
run (seed) & 0 & 1 & 2 & 3 & 4\\\midrule
20k: $\ell(\theta^\star)-\ell(\hat\theta)$ [nats] & 149.8 & 137.7 & 122.4 & 151.7 & 149.6\\
5k:\ \ $\ell(\theta^\star)-\ell(\hat\theta)$ [nats] & 168.9 & 164.3 & 116.5 & 177.6 & 170.3\\
expected, $4096\sum_n\KL_n(\theta^\star\|\hat\theta)$ (20k) & 124 & 141 & 130 & 132 & 139\\\bottomrule
\end{tabular}
\end{center}
A likelihood ratio of $e^{116}$--$e^{178}$ is decisive. The NPE posterior concentrates on parameters that the data exclude. The clean PPC does not contradict this: PPC compares \emph{summaries} within 95\% bands, a far coarser check than the likelihood.

\paragraph{3. The summaries carry the information; the network does not use it.} We simulated 300 datasets at $\theta^\star$ and 300 at $\hat\theta$ (seed 0, 20k) with the \code{moments} summary. A regularised logistic regression separates them with 5-fold CV accuracy \textbf{0.968} (0.722 at 165 clones per endpoint). The most shifted features belong to the earliest endpoint (h4: sd of $\log(1+B)$, median of $\log(1+A)$) and to the A--B correlation. In Gaussian-equivalent terms, accuracy 0.968 corresponds to an expected log-LR of $\approx7$ nats. So the summary keeps about 5\% of the exact 124 nats along this direction, and still more than enough to reject $\hat\theta$ ($\approx900{:}1$). The amortised NPE, trained on 20k draws spread over a 10-dimensional prior, does not resolve this.

\paragraph{Conclusion.} In GW the B row is not structurally non-identifiable. It is \emph{practically} non-identifiable at the designs of E0, X5a, X5c and the real-like design (Proposition~\ref{prop:cr}). At 4096 clones per endpoint it is identifiable from the data, from the exact likelihood and to a useful extent from the summaries, and the failure reported there is an inference failure (amortised NPE with uninformative coverage of the relevant region). The report's recommendations survive: fix or measure B divisions, because realistic designs cannot decide them. Their justification and wording should change (\S\ref{sec:wording}).

%======================================================================
\section{Profile: how flat is the likelihood away from $\theta^\star$?}\label{sec:profile}
The Fisher information is local. For the global picture we compute the profile
\[
K_{\mathcal D}(\delta)=\min\Big\{\sum_n w_n\KL\big(\tilde p_n(\theta^\star)\,\|\,\tilde p_n(\theta)\big):\ \theta\in\Delta_5\times\Delta_5,\ a_B+b_B+c_B=\delta\Big\},
\]
the per-clone cost of the best-fitting model in which B divides with probability $\delta$. All other 9 coordinates are optimised (L-BFGS, three starts, exact law). Then $N\,K(\delta)$ is the expected log-likelihood ratio against the best ``B divides'' explanation with $N$ organoids in total, split in the design's proportions ($w_n=N_n/N$).

\begin{table}[h]\centering\small
\begin{tabular}{@{}lrrrrrr@{}}\toprule
B-division mass $\delta$ & 0.02 & 0.05 & 0.1 & 0.2 & 0.3 & 0.5 \\\midrule
\multicolumn{7}{@{}l}{\emph{real-like} $\{6{:}165,10{:}101\}$ (weights 165/266, 101/266)}\\
$K(\delta)$ per clone & $9.8\cdot10^{-6}$ & $5.4\cdot10^{-5}$ & $1.9\cdot10^{-4}$ & $5.8\cdot10^{-4}$ & $9.3\cdot10^{-4}$ & $1.36\cdot10^{-3}$ \\
$266\,K(\delta)$ [nats] & 0.003 & 0.014 & 0.05 & 0.15 & 0.25 & 0.36 \\
best B$\to$A / A death & 0.78 / 0.04 & 0.76 / 0.05 & 0.71 / 0.07 & 0.62 / 0.10 & 0.55 / 0.12 & 0.39 / 0.16 \\\midrule
\multicolumn{7}{@{}l}{\emph{E0 at 4096} $\{4,6,10\}$, equal weights}\\
$K(\delta)$ per clone & $2.9\cdot10^{-5}$ & $1.6\cdot10^{-4}$ & $5.5\cdot10^{-4}$ & $1.6\cdot10^{-3}$ & $2.8\cdot10^{-3}$ & $4.2\cdot10^{-3}$ \\
$12\,288\,K(\delta)$ [nats] & 0.35 & 1.9 & 6.7 & 20 & 34 & 51 \\
best B$\to$A / A death & 0.78 / 0.04 & 0.76 / 0.05 & 0.72 / 0.07 & 0.64 / 0.10 & 0.54 / 0.12 & 0.33 / 0.15 \\\bottomrule
\end{tabular}
\caption{Profile KL $K(\delta)$ (truth $\delta=0$, B$\to$A $=0.8$, A death 0.03) and the B$\to$A and A death of the best-fitting model at each $\delta$. Expected log-LR of truth vs.\ that model $=$ number of organoids $\times K(\delta)$. Full minimisers: \code{scripts/b_row_theory/profile.jsonl}.}\label{tab:profile}
\end{table}

\begin{computation}[Global profile]\label{comp:profile}
(i) $K(\delta)>0$ for every $\delta$ on the grid. Near 0 it matches the Fisher prediction $\delta^2/(2\sigma_\delta^2)$, with $\sigma^2_\delta$ the per-clone CR variance of the design (E0 at 4096: predicted $2.8\cdot10^{-5}$ at $\delta=0.02$, computed $2.9\cdot10^{-5}$). Above $\delta\approx0.2$, $K$ grows sub-quadratically: explanations in which B divides a lot are relatively cheap.
(ii) \emph{Real-like design:} the best model with $\delta=0.5$ (B$\to$A $=0.39$) has an expected log-LR of only 0.36 nats against the truth, a likelihood ratio of about $1.4$. Every $\delta\le0.5$ is statistically indistinguishable from $\delta=0$ with 266 organoids. This is the non-local form of Proposition~\ref{prop:cr}. Three nats of expected evidence against $\delta=0.5$ would need $\approx2\,200$ organoids in the real-like split, or $\approx720$ split equally over $\{4,6,10\}$.
(iii) \emph{E0 at 4096:} the same alternative costs 51 nats, so it is excluded. Its parameters (B$\to$A 0.33, A death 0.15, B$\to$AA 0.34) are close to the NPE posterior mean (B$\to$A 0.32, A death 0.09, B$\to$AA 0.21). NPE found the bottom of the valley but not its depth.
(iv) \emph{The compensation} is mainly move $e_1$ (switch $\to$ divide-into-AA) in B's row, with the extra death \emph{moved to A's row}. A death rises from 0.03 to 0.16 while B death stays at 0.06--0.10. This is the precise form of ``the confounding leaks into A's row''.
\end{computation}

\noindent X5c (\{6,10\}$\times$\{32,128,512\}) is consistent with this. Using the real-like profile as an approximation (its split differs slightly), 1\,024 organoids give $\approx1.4$ nats of expected evidence against $\delta=0.5$. So its failure is mainly the likelihood's, not only the network's.

%======================================================================
\section{Suggested changes to the report's wording}\label{sec:wording}
\begin{itemize}
\item Replace ``structural non-identifiability'' (E0 at 4096, Part~I summary, ``Decided: structural'') with: \emph{``practically non-identifiable at realistic designs (Cram\'er--Rao sd of B$\to$A 0.27--0.28 > prior sd); locally identifiable in principle (positive definite exact Fisher information); at 4096 clones per endpoint the data and the summaries identify it, but amortised NPE does not.''}
\item ``A clean PPC is no evidence that the B row is right'' stays true. Add: the exact likelihood ratio against the NPE answer is $e^{120}$--$e^{180}$ on the same data.
\item ``Biological constraints help more than data'' holds at realistic sizes, with the caveat that more \emph{early} data help much more than more late data (Proposition~\ref{prop:decay}).
\item E0's first attribution (prior tail) and its revision stay as they are. The revision ``the confounding is the root cause'' is correct as a description of practical non-identifiability.
\end{itemize}

%======================================================================
\section{Scope, caveats and open questions}\label{sec:open}
\begin{itemize}
\item \textbf{Global identifiability} (injectivity) is not proved. The profile of \S\ref{sec:profile} is numerical evidence over a grid of $\delta$, from finitely many starts.
\item \textbf{One truth.} All computations are at $\theta^\star$. Theorems~\ref{thm:mean}--\ref{thm:varw} hold for every $\theta$; the Fisher numbers do not. The same scripts compute them at any other truth (e.g.\ the X5a truths).
\item \textbf{Boundary.} $\theta^\star$ has $\delta=0$ on the boundary. The MLE of $\delta$ is asymptotically a censored normal, and the CR bound is the local scale, not an exact sd.
\item \textbf{CR vs posterior.} The CR bound concerns unbiased estimators. Posteriors with informative priors can have smaller sd but are then biased, which is the E0/X5 picture. Asymptotically (data dominating) the posterior sd approaches the CR sd.
\item \textbf{Bellman--Harris.} Not analysed. Continuous lifetimes add a free time scale. The BH ridge of Part~III and the larger B$\to$A error under BH at 4096 suggest weaker information there. An exact BH analysis would need the age-dependent pgf (integral equation), which is not done here.
\item \textbf{Generations vs days.} The early-endpoint recommendation is in GW generations. Its translation to days, and to what can be counted and typed in 2--8-cell organoids, is a wet-lab question.
\item \textbf{Inference fix at 4096.} The analysis suggests likelihood-based inference (exact GW likelihood here, or sequential/local NPE around the data), or summaries targeted at early endpoints. It does not show that a particular NPE configuration would succeed.
\end{itemize}

%======================================================================
\section{A toolkit for similar situations}\label{sec:toolkit}
The analysis above follows a sequence that applies to any simulator-based model where a parameter ``is not recovered''. Each step names the readings in \S\ref{sec:reading} that teach it.
\begin{enumerate}
\item \textbf{Look for exact invariances first} (moments, symmetries). Which parameter combinations leave the mean, the variance, etc.\ unchanged? (Theorems~\ref{thm:mean}--\ref{thm:varw}; [B1, B2, B5].)
\item \textbf{Separate structural from practical.} Is the Fisher information singular (structural, at least locally), or just small? A singular information means a flat likelihood; a small one means more or better data. ([I1--I4].)
\item \textbf{Profile the likelihood} along the suspect direction, re-optimising everything else. It shows how deep the valley is, globally and not only near the truth, and which other parameters compensate. ([I2, I3].)
\item \textbf{Compare the method with the bound.} If the posterior is much wider or more biased than the Cram\'er--Rao bound, the method is losing information. Check where: exact likelihood ratio on the observed data, then a classifier two-sample test on the summaries, then the network. ([S5--S8, S11].)
\item \textbf{Use the information to design.} Fisher information per organoid and per endpoint says which measurements help; this is (local) optimal experimental design. ([D1, D2].)
\item \textbf{Treat clean diagnostics with care.} SBC averages over the prior, and a PPC on summaries is coarser than the likelihood. Neither shows identifiability. ([S6, S9, S10].)
\end{enumerate}

%======================================================================
\section{Further reading}\label{sec:reading}
Suggested order for a first pass: B3 or B4 (branching processes for biologists), I2 and I3 (identifiability in practice), S1 and S8 (simulation-based inference and its failure modes), then D2. Starred items are the most relevant to this note.

\subsection*{Branching processes (the model)}
\begin{itemize}
\item[B1$^\star$] K.\,B.\ Athreya and P.\,E.\ Ney, \emph{Branching Processes}, Springer, 1972 (Dover reprint 2004). Ch.~V: multitype processes, the mean matrix, moment recursions and the Kesten--Stigum limit, i.e.\ the background of Theorems~\ref{thm:second} and~\ref{thm:varw}.
\item[B2] T.\,E.\ Harris, \emph{The Theory of Branching Processes}, Springer, 1963. The classical source for probability generating functions and the covariance recursion~\eqref{eq:cov}; also covers age-dependent (Bellman--Harris) processes.
\item[B3$^\star$] M.\ Kimmel and D.\,E.\ Axelrod, \emph{Branching Processes in Biology}, 2nd ed., Springer, 2015. The most accessible entry point; cell proliferation, multitype and age-dependent models, with biological examples.
\item[B4] P.\ Haccou, P.\ Jagers and V.\,A.\ Vatutin, \emph{Branching Processes: Variation, Growth, and Extinction of Populations}, Cambridge University Press, 2005. Population-biology viewpoint; good on multitype growth and extinction.
\item[B5] H.\ Kesten and B.\,P.\ Stigum, ``A limit theorem for multidimensional Galton--Watson processes'', \emph{Ann.\ Math.\ Statist.}\ 37 (1966) 1211--1223. Origin of $Z_n\approx\rho^nWu$.
\item[B6] P.\ Guttorp, \emph{Statistical Inference for Branching Processes}, Wiley, 1991. Likelihood and estimation for branching processes, including what can and cannot be estimated from generation counts.
\end{itemize}

\subsection*{Clonal dynamics in stem-cell biology (the application)}
\begin{itemize}
\item[C1$^\star$] H.\,J.\ Snippert et al., ``Intestinal crypt homeostasis results from neutral competition between symmetrically dividing Lgr5 stem cells'', \emph{Cell} 143 (2010) 134--144. LGR5$^+$ clone-size data interpreted with stochastic fate models.
\item[C2] E.\ Clayton et al., ``A single type of progenitor cell maintains normal epidermis'', \emph{Nature} 446 (2007) 185--189. A classic inference of division and fate probabilities from clone-size distributions.
\item[C3] A.\,M.\ Klein and B.\,D.\ Simons, ``Universal patterns of stem cell fate in cycling adult tissues'', \emph{Development} 138 (2011) 3103--3111. Why clone-size distributions often fix only some combinations of fate parameters.
\item[C4] S.\ Rulands and B.\,D.\ Simons, ``Tracing cellular dynamics in tissue development, maintenance and disease'', \emph{Curr.\ Opin.\ Cell Biol.}\ 43 (2016) 38--45. Short review of lineage tracing and its quantitative analysis.
\end{itemize}

\subsection*{Identifiability, Fisher information, profile likelihood}
\begin{itemize}
\item[I1] T.\,J.\ Rothenberg, ``Identification in parametric models'', \emph{Econometrica} 39 (1971) 577--591. Local identifiability $\Leftrightarrow$ non-singular Fisher information (under regularity), the basis of Computation~\ref{comp:fim}.
\item[I2$^\star$] A.\ Raue et al., ``Structural and practical identifiability analysis of partially observed dynamical models by exploiting the profile likelihood'', \emph{Bioinformatics} 25 (2009) 1923--1929. The structural vs practical distinction and the profile-likelihood method used in \S\ref{sec:profile}.
\item[I3$^\star$] F.-G.\ Wieland, A.\,L.\ Hauber, M.\ Rosenblatt, C.\ T\"onsing and J.\ Timmer, ``On structural and practical identifiability'', \emph{Curr.\ Opin.\ Syst.\ Biol.}\ 25 (2021) 60--69. A recent, readable review of the same distinction, with pitfalls.
\item[I4] R.\,N.\ Gutenkunst et al., ``Universally sloppy parameter sensitivities in systems biology models'', \emph{PLoS Comput.\ Biol.}\ 3 (2007) e189; and M.\,K.\ Transtrum et al., ``Perspective: Sloppiness and emergent theories in physics, biology, and beyond'', \emph{J.\ Chem.\ Phys.}\ 143 (2015) 010901. Why Fisher eigenvalues spread over many decades (cf.\ Table~\ref{tab:fim}) and what that means.
\item[I5] E.\,L.\ Lehmann and G.\ Casella, \emph{Theory of Point Estimation}, 2nd ed., Springer, 1998. Fisher information and the Cram\'er--Rao bound, with proofs.
\item[I6] A.\,W.\ van der Vaart, \emph{Asymptotic Statistics}, Cambridge University Press, 1998. Asymptotics of the MLE and of posteriors (Bernstein--von Mises): why the CR bound is the right yardstick for large samples.
\item[I7] S.\,G.\ Self and K.-Y.\ Liang, ``Asymptotic properties of maximum likelihood estimators and likelihood ratio tests under nonstandard conditions'', \emph{J.\ Amer.\ Statist.\ Assoc.}\ 82 (1987) 605--610. Parameters on the boundary (here $\delta=0$).
\end{itemize}

\subsection*{Simulation-based inference: methods, summaries, diagnostics}
\begin{itemize}
\item[S1$^\star$] K.\ Cranmer, J.\ Brehmer and G.\ Louppe, ``The frontier of simulation-based inference'', \emph{PNAS} 117 (2020) 30055--30062. Overview of NPE, NLE, NRE and where information is lost.
\item[S2] G.\ Papamakarios and I.\ Murray, ``Fast $\epsilon$-free inference of simulation models with Bayesian conditional density estimation'', NeurIPS 2016; D.\,S.\ Greenberg, M.\ Nonnenmacher and J.\,H.\ Macke, ``Automatic posterior transformation for likelihood-free inference'', ICML 2019. Neural posterior estimation, amortised and sequential.
\item[S3] G.\ Papamakarios, D.\ Sterratt and I.\ Murray, ``Sequential neural likelihood'', AISTATS 2019. Learning the likelihood instead of the posterior; an alternative when the posterior is hard to amortise.
\item[S4] M.\ Deistler, P.\,J.\ Gon\c{c}alves and J.\,H.\ Macke, ``Truncated proposals for scalable and hassle-free simulation-based inference'', NeurIPS 2022. Concentrating simulations near the data: the remedy suggested by \S\ref{sec:e0large}.
\item[S5] J.-M.\ Lueckmann et al., ``Benchmarking simulation-based inference'', AISTATS 2021. How methods compare against reference posteriors, the logic of item~4 of \S\ref{sec:toolkit}.
\item[S6$^\star$] J.\ Hermans et al., ``A trust crisis in simulation-based inference? Your posterior approximations can be unfaithful'', \emph{TMLR} 2022. Over-confident or biased NPE posteriors, as at 4096 clones.
\item[S7] P.\ Fearnhead and D.\ Prangle, ``Constructing summary statistics for approximate Bayesian computation: semi-automatic approximate Bayesian computation'', \emph{J.\ R.\ Stat.\ Soc.\ B} 74 (2012) 419--474; Y.\ Chen et al., ``Neural approximate sufficient statistics for implicit models'', ICLR 2021. Choosing or learning summaries that keep the information (cf.\ the 5\% kept by \code{moments}).
\item[S8$^\star$] S.\,A.\ Sisson, Y.\ Fan and M.\ Beaumont (eds.), \emph{Handbook of Approximate Bayesian Computation}, CRC Press, 2018. Broad reference on likelihood-free inference, summary statistics and diagnostics.
\item[S9] S.\ Talts, M.\ Betancourt, D.\ Simpson, A.\ Vehtari and A.\ Gelman, ``Validating Bayesian inference algorithms with simulation-based calibration'', arXiv:1804.06788 (2018). SBC and what it does \emph{not} check.
\item[S10] A.\ Gelman et al., \emph{Bayesian Data Analysis}, 3rd ed., CRC Press, 2013, Ch.~6. Posterior predictive checks and their limits.
\item[S11] D.\ Lopez-Paz and M.\ Oquab, ``Revisiting classifier two-sample tests'', ICLR 2017. The C2ST used in E0 and the two-point test of \S\ref{sec:e0large}.
\item[S12] S.\,N.\ Wood, ``Statistical inference for noisy nonlinear ecological systems'', \emph{Nature} 466 (2010) 1102--1104. Synthetic likelihood: a Gaussian model of the summaries. Note its bias with noisy finite-difference derivatives, which made us drop such an estimate here.
\item[S13] A.\ Tejero-Cantero et al., ``sbi: A toolkit for simulation-based inference'', \emph{J.\ Open Source Softw.}\ 5 (2020) 2505. The library used by this project.
\end{itemize}

\subsection*{Experimental design}
\begin{itemize}
\item[D1] A.\,C.\ Atkinson, A.\,N.\ Donev and R.\,D.\ Tobias, \emph{Optimum Experimental Designs, with SAS}, Oxford University Press, 2007. Fisher-information-based (D-, A-, c-optimal) design; the c-optimal criterion for one parameter is exactly the CR variance of Table~\ref{tab:cr}.
\item[D2$^\star$] K.\ Chaloner and I.\ Verdinelli, ``Bayesian experimental design: a review'', \emph{Statist.\ Sci.}\ 10 (1995) 273--304. Design when the prior matters, as it does at realistic sizes here.
\end{itemize}

\raggedright
\paragraph{Reproducibility.} \code{scripts/b_row_theory/}: \code{gwexact.py} (exact law, derivatives, Fisher information, KL), \code{fisher.py}, \code{fisher2.py} (Tables~\ref{tab:fim}, \ref{tab:cr}), \code{llr.py} (\S\ref{sec:e0large}, item 2; reads \code{results/e4/gw_xxl_two_type_*}), \code{summary_twopoint.py} (item 3), \code{varW.py} (Theorem~\ref{thm:varw} check), \code{profile.py} (\S\ref{sec:profile}). Run with \code{uv run python scripts/b_row_theory/<script>.py} from the repository root.

\end{document}
